%% =====================================================================
%% Trabalho 06 - Aproximação Funcional com MLP
%% EL056 - Redes Neurais Artificiais - PPGEELT/UFU
%% Compilar com pdfLaTeX + BibTeX (Overleaf: Menu > Compiler > pdfLaTeX)
%% =====================================================================
\documentclass[12pt,oneside,a4paper,english,brazil]{abntex2}

\usepackage[T1]{fontenc}
\usepackage[utf8]{inputenc}
\usepackage{lmodern}
\usepackage{microtype}
\usepackage{amsmath,amssymb,bm}
\usepackage{graphicx}
\usepackage{booktabs}
\usepackage{float}
\usepackage{subcaption}
\usepackage[portuguese,ruled,vlined,linesnumbered]{algorithm2e}
\usepackage{xcolor}
\usepackage[alf]{abntex2cite}

\graphicspath{{figuras/}}

%% ---------------------------------------------------------------------
%% Dados da capa e folha de rosto
%% ---------------------------------------------------------------------
\titulo{Trabalho 06: Aproximação Funcional com Perceptron Multicamadas}
\autor{Lucas A. Martins}
\local{Uberlândia -- MG}
\data{2026}
\instituicao{%
  Universidade Federal de Uberlândia -- UFU
  \par
  Faculdade de Engenharia Elétrica -- FEELT
  \par
  Programa de Pós-Graduação em Engenharia Elétrica -- PPGEELT}
\tipotrabalho{Relatório técnico}
\preambulo{Relatório apresentado à disciplina EL056 -- Redes Neurais Artificiais,
do Programa de Pós-Graduação em Engenharia Elétrica da Universidade Federal de
Uberlândia, como requisito parcial para aprovação.\\[1em]
Matrícula: 12622EEL006}
\orientador[Professor:]{Prof. Dr. Keiji Yamanaka}

\makeatletter
\hypersetup{
  pdftitle={\@title}, pdfauthor={\@author},
  pdfsubject={EL056 - Redes Neurais Artificiais},
  colorlinks=true, linkcolor=blue!50!black, citecolor=blue!50!black,
  urlcolor=blue!50!black, bookmarksdepth=4}
\makeatother

\setlength{\parindent}{1.25cm}
\setlength{\parskip}{0.1cm}
\SetAlgoCaptionSeparator{ --}
\newcommand{\mse}{\mathrm{MSE}}

%% Relatório curto: seções numeradas sem capítulo e numeração contínua
\renewcommand{\thesection}{\arabic{section}}
\counterwithout{figure}{chapter}
\counterwithout{table}{chapter}
\counterwithout{equation}{chapter}

\begin{document}
\selectlanguage{brazil}
\frenchspacing

\imprimircapa
\imprimirfolhaderosto

\textual

%% =====================================================================
\section{Introdução}
%% =====================================================================
Aproximar uma função desconhecida $f:\mathbb{R}\to\mathbb{R}$ a partir de um
conjunto finito de amostras $\{(x_p,t_p)\}_{p=1}^{P}$ é um dos problemas
clássicos de redes neurais: a rede deve produzir uma função contínua que
passe próxima dos pontos observados e que tenha comportamento razoável entre
eles. Neste trabalho, os $P=11$ pontos fornecidos no enunciado
(Tabela~\ref{tab:dados}) são aproximados por um Perceptron Multicamadas (MLP)
treinado por retropropagação do erro (\textit{backpropagation}), implementado
integralmente em NumPy.

\begin{table}[H]
\centering
\caption{Pontos amostrados fornecidos no enunciado.}
\label{tab:dados}
\footnotesize\setlength{\tabcolsep}{2.5pt}
\begin{tabular}{l*{11}{r}}
\toprule
$x$ & 0 & 0,1 & 0,2 & 0,3 & 0,4 & 0,5 & 0,6 & 0,7 & 0,8 & 0,9 & 1,0\\
$t$ & $-$0,9602 & $-$0,5770 & $-$0,0729 & 0,3771 & 0,6405 & 0,6600 & 0,4609 & 0,1336 & $-$0,2013 & $-$0,4344 & $-$0,5000\\
\bottomrule
\end{tabular}
\legend{Fonte: enunciado do Trabalho 06.}
\end{table}

Além do treinamento da rede pedida, investigam-se o efeito do número de
neurônios ocultos, da taxa de aprendizado, do termo de momento, da
inicialização dos pesos e do modo de treinamento, bem como a capacidade de
generalização e de extrapolação da rede obtida.

%% =====================================================================
\section{Fundamentação teórica}
%% =====================================================================

\subsection{O Perceptron Multicamadas}
Um neurônio artificial calcula a soma ponderada de suas entradas acrescida de
um bias, $\mathit{net}=\mathbf{w}^{\top}\mathbf{x}+b$, e aplica uma função de
ativação $\varphi(\mathit{net})$~\cite{haykin2009}. Um MLP organiza esses
neurônios em camadas totalmente conectadas. Para o problema escalar deste
trabalho usa-se a topologia $1$--$N$--$1$: uma entrada, $N$ neurônios ocultos
com tangente hiperbólica e um neurônio de saída linear. Em notação matricial,
para um lote $\mathbf{X}\in\mathbb{R}^{1\times P}$:
\begin{equation}
\mathbf{net}^{(1)}=\mathbf{W}_1\mathbf{X}+\mathbf{b}_1,\qquad
\mathbf{H}=\tanh\!\big(\mathbf{net}^{(1)}\big),\qquad
\mathbf{Y}=\mathbf{W}_2\mathbf{H}+b_2,
\label{eq:forward}
\end{equation}
com $\mathbf{W}_1,\mathbf{b}_1\in\mathbb{R}^{N\times1}$,
$\mathbf{W}_2\in\mathbb{R}^{1\times N}$ e $b_2\in\mathbb{R}$, totalizando
$3N+1$ parâmetros livres.

\subsection{Funções de ativação}
A tangente hiperbólica $\varphi(v)=\tanh(v)$ é sigmoidal, limitada em
$(-1,1)$, antissimétrica e diferenciável, com derivada expressa pela própria
saída:
\begin{equation}
\varphi'(v)=1-\tanh^2(v)=1-h^2.
\end{equation}
A antissimetria acelera o treinamento em relação à sigmoide
logística~\cite{haykin2009}. Na saída adota-se a ativação linear
$\varphi(v)=v$, com $\varphi'(v)=1$. Assim a amplitude da resposta não fica
limitada a $(-1,1)$ e a rede implementa uma combinação linear de funções
sigmoidais deslocadas e escaladas.

\subsection{Teorema da Aproximação Universal}
\citeonline{cybenko1989} demonstrou que somas finitas da forma
\begin{equation}
G(x)=\sum_{j=1}^{N}\alpha_j\,\sigma\!\left(w_j x+\theta_j\right)
\label{eq:uat}
\end{equation}
com $\sigma$ sigmoidal contínua são densas no espaço das funções contínuas
sobre um compacto: para toda $f$ contínua e todo $\varepsilon>0$ existe $N$
finito tal que $|G(x)-f(x)|<\varepsilon$ em todo o domínio.
\citeonline{hornik1991} mostrou que a propriedade decorre da própria
arquitetura em camadas, valendo para qualquer ativação limitada e não
constante. A Eq.~\eqref{eq:uat} é exatamente a rede $1$--$N$--$1$ da
Eq.~\eqref{eq:forward}.

O teorema é de \emph{existência}. Ele não informa quantos neurônios são
necessários, não garante que o gradiente descendente encontre os pesos
adequados, não trata da generalização a partir de poucas amostras e não diz
nada sobre o comportamento da rede fora do domínio compacto considerado.
Esses aspectos são avaliados empiricamente na Seção~\ref{cap:resultados}.

\subsection{Algoritmo de retropropagação}
O treinamento minimiza o erro quadrático
\begin{equation}
E=\frac{1}{2P}\sum_{p=1}^{P}\left(t_p-y_p\right)^2,
\end{equation}
cujo gradiente é obtido pela regra da cadeia, propagando o erro da saída para
a entrada~\cite{rumelhart1986}. Definem-se os gradientes locais
$\delta=\partial E/\partial\mathit{net}$ (a menos do fator $1/P$). Na saída
linear,
\begin{equation}
\bm{\delta}^{(2)}=\mathbf{Y}-\mathbf{T},
\end{equation}
e, na camada oculta, o erro é retropropagado pelos pesos $\mathbf{W}_2$ e
modulado pela derivada da ativação:
\begin{equation}
\bm{\delta}^{(1)}=\left(\mathbf{W}_2^{\top}\bm{\delta}^{(2)}\right)\odot\left(1-\mathbf{H}^{2}\right).
\end{equation}
Os gradientes em relação aos parâmetros são o produto do $\delta$ de cada
neurônio pelo sinal que chega a ele:
\begin{equation}
\frac{\partial E}{\partial\mathbf{W}_2}=\frac1P\bm{\delta}^{(2)}\mathbf{H}^{\top},\quad
\frac{\partial E}{\partial b_2}=\frac1P\sum_p\delta^{(2)}_p,\quad
\frac{\partial E}{\partial\mathbf{W}_1}=\frac1P\bm{\delta}^{(1)}\mathbf{X}^{\top},\quad
\frac{\partial E}{\partial\mathbf{b}_1}=\frac1P\sum_p\bm{\delta}^{(1)}_{:,p}.
\end{equation}
A implementação foi validada comparando esses gradientes com diferenças
finitas centradas. O erro relativo foi inferior a $10^{-8}$ em todos os
parâmetros.

\subsection{Gradiente descendente com momento}
Os pesos são atualizados na direção oposta ao gradiente, com taxa de
aprendizado $\alpha$. O termo de momento $\beta$ acumula uma fração da
atualização anterior~\cite{rumelhart1986,haykin2009}:
\begin{equation}
\Delta\mathbf{w}(k)=-\alpha\,\nabla_{\mathbf{w}}E+\beta\,\Delta\mathbf{w}(k-1),\qquad
\mathbf{w}(k+1)=\mathbf{w}(k)+\Delta\mathbf{w}(k).
\end{equation}
O momento funciona como um filtro passa-baixas sobre as atualizações: acelera
o avanço em direções de gradiente consistente (platôs) e amortece oscilações
em vales estreitos. No modo \emph{em lote} faz-se uma atualização por época
com o gradiente médio. No modo \emph{online} (padrão a padrão) faz-se uma
atualização por amostra, em ordem aleatória a cada época. O procedimento é
resumido no Algoritmo~\ref{alg:bp}.

\begin{algorithm}[H]
\caption{Treinamento da MLP $1$--$N$--$1$ por retropropagação com momento (modo em lote).}
\label{alg:bp}
\KwIn{amostras $(\mathbf{X},\mathbf{T})$; $N$, $\alpha$, $\beta$, $\mathit{tol}$, $\mathit{max\_epocas}$}
Inicializar $\mathbf{W}_1,\mathbf{b}_1,\mathbf{W}_2,b_2$ (Nguyen--Widrow) e velocidades $\Delta\mathbf{w}\gets\mathbf{0}$\;
\For{$k\gets1$ \KwTo $\mathit{max\_epocas}$}{
  $\mathbf{H}\gets\tanh(\mathbf{W}_1\mathbf{X}+\mathbf{b}_1)$;\quad $\mathbf{Y}\gets\mathbf{W}_2\mathbf{H}+b_2$\tcp*{direta}
  $\bm{\delta}^{(2)}\gets\mathbf{Y}-\mathbf{T}$;\quad $\bm{\delta}^{(1)}\gets(\mathbf{W}_2^{\top}\bm{\delta}^{(2)})\odot(1-\mathbf{H}^2)$\tcp*{retropropag.}
  \ForEach{parâmetro $\mathbf{w}$ com gradiente $\mathbf{g}$}{
    $\Delta\mathbf{w}\gets\beta\,\Delta\mathbf{w}-\alpha\,\mathbf{g}$;\quad $\mathbf{w}\gets\mathbf{w}+\Delta\mathbf{w}$\;
  }
  \lIf{$\mse=\frac1P\sum_p(t_p-y_p)^2<\mathit{tol}$}{\textbf{parar}}
}
\end{algorithm}

\subsection{Inicialização de Nguyen--Widrow}
Com pesos pequenos e aleatórios, as regiões de maior inclinação das tanh
(onde $|\mathit{net}|$ é pequeno) tendem a se concentrar em uma mesma parte do
domínio, e vários neurônios acabam fazendo o mesmo papel.
\citeonline{nguyen1990} propuseram redimensionar os pesos da camada oculta para
que essas regiões ativas se distribuam ao longo do intervalo de entrada. Com
$n$ entradas e $N$ neurônios ocultos, define-se $\beta_{\mathrm{NW}}=0{,}7\,N^{1/n}$,
sorteia-se $\mathbf{w}_j\sim\mathcal{U}(-0{,}5;\,0{,}5)$, normaliza-se cada vetor
para $\|\mathbf{w}_j\|=\beta_{\mathrm{NW}}$ e sorteiam-se os bias em
$\mathcal{U}(-\beta_{\mathrm{NW}};\,\beta_{\mathrm{NW}})$~\cite{fausett1994}.

\subsection{Viés, variância e generalização}
Uma rede com poucos neurônios tem viés alto e não consegue representar a
forma da função (subajuste). Uma rede com muitos parâmetros em relação ao
número de amostras tem variância alta e pode interpolar os pontos oscilando
entre eles (sobreajuste). Com apenas 11 amostras, o erro de treino não serve
para escolher $N$, porque ele sempre diminui com a capacidade da rede. Por
isso a generalização foi estimada por validação cruzada
\emph{leave-one-out} (LOO): para cada ponto $i$, treina-se a rede com os
outros 10 pontos e mede-se o erro quadrático em $x_i$. O MSE-LOO é a média
desses 11 erros.

%% =====================================================================
\section{Metodologia}
%% =====================================================================
A rede, a retropropagação e o otimizador foram implementados do zero em
Python/NumPy. O programa está no arquivo \texttt{trabalho06\_mlp.py}, e a
demonstração passo a passo, no \textit{notebook} \texttt{trabalho06\_mlp.ipynb}. A configuração base foi: MLP
$1$--$5$--$1$, tanh na camada oculta, saída linear, inicialização de
Nguyen--Widrow, treinamento em lote, $\alpha=0{,}05$, $\beta=0{,}9$, parada em
$\mse<10^{-3}$ ou após 20\,000 épocas, e semente 42. Os experimentos foram os
seguintes:
(i)~treinamento base;
(ii)~$N\in\{1,2,3,5,10,20\}$ treinados por 50\,000 épocas, sem parada
antecipada;
(iii)~$\alpha\in\{0{,}01;\,0{,}05;\,0{,}1;\,0{,}3\}$ combinado com
$\beta\in\{0;\,0{,}5;\,0{,}9\}$,
medindo as épocas até $\mse<10^{-3}$ (mediana de 10 sementes);
(iv)~inicialização uniforme $\times$ Nguyen--Widrow e modo em lote $\times$
online (10 sementes);
(v)~MSE-LOO para cada $N$; e
(vi)~avaliação da rede em $x\in[-1,1]$ para analisar a extrapolação.

%% =====================================================================
\section{Resultados e discussão}\label{cap:resultados}
%% =====================================================================

\subsection{Treinamento base}
A rede $1$--$5$--$1$ atingiu $\mse=9{,}99\times10^{-4}$ em \textbf{89 épocas}.
A Figura~\ref{fig:base} reproduz o formato das figuras do enunciado. A curva
aprendida é suave, tem pico próximo de $x\approx0{,}46$ e passa próxima dos 11
pontos. O maior desvio ocorre em $x=1{,}0$ ($y=-0{,}572$ para $t=-0{,}500$).
Esse ajuste é mais fiel aos pontos do que a curva ilustrativa do enunciado,
cujo pico fica perto de $0{,}5$ e que se afasta visivelmente dos pontos em
$x=0$ e $x\in\{0{,}4;\,0{,}5\}$. Continuando o treinamento sem parada
antecipada, o MSE cai para $8{,}2\times10^{-5}$ em 50\,000 épocas
(Figura~\ref{fig:curva}), mas quase toda a redução acontece nas primeiras
centenas de épocas. Os pesos finais estão na Tabela~\ref{tab:pesos}.

\begin{figure}[H]
\centering
\begin{subfigure}{0.38\textwidth}\centering
  \includegraphics[width=\linewidth]{dados.png}
  \caption{Pontos amostrados.}
\end{subfigure}\hfill
\begin{subfigure}{0.38\textwidth}\centering
  \includegraphics[width=\linewidth]{aproximacao_base.png}
  \caption{Aproximação pela MLP $1$--$5$--$1$.}
\end{subfigure}
\caption{Dados do enunciado e função aproximada pela rede treinada.}
\label{fig:base}
\legend{Fonte: elaborado pelo autor.}
\end{figure}

\begin{figure}[H]
\centering
\includegraphics[width=0.62\textwidth]{curva_aprendizado.png}
\caption{Curva de aprendizado ($\mse\times$ época, escalas logarítmicas).}
\label{fig:curva}
\legend{Fonte: elaborado pelo autor.}
\end{figure}

\begin{table}[H]
\centering
\caption{Pesos finais da rede base ($N=5$, 89 épocas).}
\label{tab:pesos}
\small
\begin{tabular}{crrrrr}
\toprule
Neurônio oculto $j$ & 1 & 2 & 3 & 4 & 5\\
\midrule
$w^{(1)}_j$ & 3,7694 & $-$3,8397 & 3,4926 & 3,5000 & $-$3,5034\\
$b^{(1)}_j$ & $-$0,6305 & 2,6729 & 1,3174 & 2,2503 & $-$0,3970\\
$w^{(2)}_j$ & 1,3552 & 0,9670 & $-$0,7331 & $-$0,2684 & $-$0,2506\\
\midrule
\multicolumn{6}{l}{Bias da saída: $b_2=-0{,}3750$}\\
\bottomrule
\end{tabular}
\legend{Fonte: elaborado pelo autor.}
\end{table}

\subsection{Número de neurônios e generalização}
A Figura~\ref{fig:neuronios} e a Tabela~\ref{tab:neuronios} mostram o efeito
de $N$. Com $N=1$ a rede só consegue representar um degrau sigmoidal (viés
alto, MSE de $0{,}14$). Dois neurônios já capturam o formato de sino, porque
a diferença de duas sigmoides deslocadas forma uma ``lombada''. O erro de
treino cai de forma monótona com $N$, e com $N=20$ (61 parâmetros para 11
pontos) a rede interpola os dados até a precisão de máquina. O MSE-LOO
(Figura~\ref{fig:loo}), por outro lado, é mínimo em $N=5$
($8{,}4\times10^{-4}$) e piora 30 vezes em $N=20$
($2{,}7\times10^{-2}$). Nesse caso, a rede passa exatamente pelos pontos
de treino, mas erra os pontos retirados. Esse é o dilema viés--variância e
justifica a escolha de $N=5$.

\begin{figure}[H]
\centering
\includegraphics[width=0.78\textwidth]{neuronios.png}
\caption{Aproximação obtida para diferentes números de neurônios ocultos (50\,000 épocas).}
\label{fig:neuronios}
\legend{Fonte: elaborado pelo autor.}
\end{figure}

\begin{table}[H]
\centering
\caption{Erro de treino e erro de validação \emph{leave-one-out} em função de $N$.}
\label{tab:neuronios}
\small
\begin{tabular}{rrrrr}
\toprule
$N$ & Parâmetros & $\mse$ treino & $\mse$ LOO & Maior $|e|$ LOO\\
\midrule
1  & 4  & $1{,}41\times10^{-1}$ & $3{,}31\times10^{-1}$ & 1,210\\
2  & 7  & $2{,}27\times10^{-4}$ & $2{,}61\times10^{-3}$ & 0,118\\
3  & 10 & $1{,}81\times10^{-4}$ & $2{,}86\times10^{-3}$ & 0,114\\
\textbf{5}  & \textbf{16} & $\bm{8{,}22\times10^{-5}}$ & $\bm{8{,}43\times10^{-4}}$ & \textbf{0,076}\\
10 & 31 & $1{,}56\times10^{-5}$ & $1{,}31\times10^{-3}$ & 0,082\\
20 & 61 & $\approx 0$ ($6{,}5\times10^{-31}$) & $2{,}68\times10^{-2}$ & 0,350\\
\bottomrule
\end{tabular}
\legend{Fonte: elaborado pelo autor.}
\end{table}

\begin{figure}[H]
\centering
\begin{subfigure}[b]{0.52\textwidth}\centering
  \includegraphics[width=\linewidth]{loo.png}
  \caption{Erro de treino e erro \emph{leave-one-out} em função de $N$.}
  \label{fig:loo}
\end{subfigure}\hfill
\begin{subfigure}[b]{0.42\textwidth}\centering
  \includegraphics[width=\linewidth]{extrapolacao.png}
  \caption{Redes treinadas avaliadas em $x\in[-1,1]$.}
  \label{fig:extrap}
\end{subfigure}
\caption{Generalização (validação cruzada) e extrapolação.}
\label{fig:gen}
\legend{Fonte: elaborado pelo autor.}
\end{figure}



\subsection{Taxa de aprendizado, momento, inicialização e modo de treinamento}
A Tabela~\ref{tab:lr} mostra que o momento acelera a convergência em todas as
taxas testadas. Com $\alpha=0{,}01$, passar de $\beta=0$ para $\beta=0{,}9$
reduz a mediana de 6\,732 para 747 épocas (cerca de 9 vezes). Isso é
coerente com a interpretação do momento como um aumento da taxa efetiva, que
se aproxima de $\alpha/(1-\beta)$ em regiões de gradiente estável. Nenhuma
combinação divergiu nesta faixa, porque o problema é pequeno e bem
condicionado. As falhas registradas com $\beta\le0{,}5$ e $\alpha\le0{,}05$
correspondem a uma semente que ficou presa em um platô e não atingiu a
tolerância em 50\,000 épocas.

\begin{table}[H]
\centering
\begin{minipage}[t]{0.56\textwidth}\centering
\caption{Épocas até $\mse<10^{-3}$ (mediana de 10 sementes; entre parênteses, quantas convergiram).}
\label{tab:lr}
\small
\begin{tabular}{rccc}
\toprule
$\alpha$ & $\beta=0$ & $\beta=0{,}5$ & $\beta=0{,}9$\\
\midrule
0,01 & 6\,732 (9/10) & 3\,367 (9/10) & 747 (10/10)\\
0,05 & 1\,347 (9/10) & 737 (10/10) & 157 (10/10)\\
0,10 & 736 (10/10) & 369 (10/10) & 117 (10/10)\\
0,30 & 246 (10/10) & 124 (10/10) & 80 (10/10)\\
\bottomrule
\end{tabular}
\legend{Fonte: elaborado pelo autor.}
\end{minipage}\hfill
\begin{minipage}[t]{0.40\textwidth}\centering
\caption{Inicialização e modo de treinamento (épocas até $\mse<10^{-3}$, mediana de 10 sementes).}
\label{tab:init}
\small\setlength{\tabcolsep}{4pt}
\begin{tabular}{lrr}
\toprule
Inicialização & Em lote & Online\\
\midrule
Uniforme & 6\,297 & 921\\
Nguyen--Widrow & \textbf{157} & \textbf{25}\\
\bottomrule
\end{tabular}
\legend{Fonte: elaborado pelo autor.}
\end{minipage}
\end{table}

A inicialização teve efeito ainda maior (Tabela~\ref{tab:init}). Com pesos
uniformes em $(-0{,}5;\,0{,}5)$ e entradas em $[0,1]$, todas as tanh operam
quase na região linear, e o treinamento em lote precisou de mediana de 6\,297
épocas. Com Nguyen--Widrow foram 157 épocas, cerca de 40 vezes menos. No modo
online a rede faz 11 atualizações por época, o que reduz o número de épocas.
Comparando pelo número de atualizações de pesos, porém, a diferença é pequena
com Nguyen--Widrow ($24{,}5\times11\approx270$ contra 157).



\subsection{Extrapolação}
A Figura~\ref{fig:extrap} avalia redes com $N=3$, $5$ e $20$ no intervalo
$[-1,1]$ dos eixos do enunciado. Dentro de $[0,1]$ (faixa amarela) as três
curvas praticamente coincidem. Fora desse intervalo elas divergem, e cada uma
segue a saturação das suas próprias tanh. A rede com $N=20$ ainda apresenta
ondulações sem nenhuma relação com os dados. Esse comportamento ilustra o
limite do Teorema da Aproximação Universal: a garantia vale apenas no domínio
compacto coberto pelas amostras, e a rede não aprende a lei que gerou os
dados.



%% =====================================================================
\section{Conclusão}
%% =====================================================================
Uma MLP $1$--$5$--$1$ com tangente hiperbólica na camada oculta e saída
linear, treinada por retropropagação com momento, aproximou os 11 pontos do
enunciado com $\mse<10^{-3}$ em 89 épocas. A curva obtida é suave e fiel aos
dados, como pedido no trabalho. Os experimentos confirmaram a teoria:
(i)~o número de neurônios deve ser escolhido pelo erro de generalização e
não pelo erro de treino, e a validação \emph{leave-one-out} indicou $N=5$ e
mostrou sobreajuste com $N=20$;
(ii)~o momento e a inicialização de Nguyen--Widrow reduzem o tempo de
treinamento em uma e duas ordens de grandeza, respectivamente; e
(iii)~a aproximação só é confiável dentro do intervalo amostrado. Como
extensões possíveis, ficam técnicas de regularização (\textit{weight decay},
parada antecipada por validação) e otimizadores de segunda ordem, como
Levenberg--Marquardt.

\postextual
\bibliography{referencias}

\end{document}
